% GENERATED FILE. Edit canonical lessons, then run npm run generate:books.
% canonical-source-hash: fnv1a64:7934efcd54f75801
% canonical-lessons: LA-C465-modius, LA-C465-sextarius, LA-C465-pondo, LA-C465-quadrans, LA-C465-nummularius

\chapter{Peck, Pint, By weight, Quarter-as, Money-changer}
\label{ch:la-peck-pint-by-weight-quarter-as-money-changer}
\hlchaptermodality{465}

\begin{chapteropening}
\textbf{By the end of this chapter:} \emph{I can talk about measures and money: a peck, a pint, things sold by weight, a small coin, and the money-changer.}

Complete the last lesson of chapter 465: I can talk about measures and money: a peck, a pint, things sold by weight, a small coin, and the money-changer.
\end{chapteropening}

\section[modius, modiī]{modius, modiī — a peck (a measure of grain)}

\label{lesson:LA-C465-modius}



\begin{quote}
Before the new one: say the Latin for to weigh out, then the Latin for to weigh carefully.
\end{quote}

\subsection*{You'll want to know: modius, modiī}
\textbf{modius, modiī} — \textquotedblleft{}a peck (a measure of grain)\textquotedblright{}.

\begin{quote}
Agricola trēs modiōs frūmentī vendit.
\end{quote}

\begin{quote}
\emph{The farmer sells three pecks of grain.}
\end{quote}

Second declension, masculine: \emph{modius, modiī m.}, genitive plural often \emph{modium}. About nine litres, the standard grain measure.

\begin{cousinweb}
\emph{modius} is built on \emph{modus} 'measure' with the suffix \emph{-ius} (EtymDB; Lewis \& Short).
\end{cousinweb}

\begin{grammarlens}[title={What you've built}]
One more word for telling, explaining and describing at B1.
\end{grammarlens}

\subsection*{Your turn}
\begin{itemize}
\raggedright
  \item \emph{Say it:} \emph{modius, modiī}
  \item \emph{Say it:} \emph{modius, modiī}, once more
  \item \emph{Say it:} a sentence of your own with \emph{modius, modiī}
  \item \emph{Recall:} say the Latin for to weigh out, then the Latin for to weigh carefully, then say \emph{modius, modiī} again
\end{itemize}

\begin{tcolorbox}[breakable,colback=teal!4,colframe=teal!35!black,title={Before you move on}]
What does modius, modiī mean? (\textquotedblleft{}A peck (a measure of grain)\textquotedblright{}.) Say it once more.
\end{tcolorbox}
\section[sextārius, sextāriī]{sextārius, sextāriī — a pint (a sixth of a larger measure)}

\label{lesson:LA-C465-sextarius}



\begin{quote}
Before the new one: say the Latin for to weigh carefully, then the Latin for a peck.
\end{quote}

\subsection*{You'll want to know: sextārius, sextāriī}
\textbf{sextārius, sextāriī} — \textquotedblleft{}a pint (a sixth of a larger measure)\textquotedblright{}.

\begin{quote}
Caupō sextārium vīnī prō duōbus assibus vendit.
\end{quote}

\begin{quote}
\emph{The innkeeper sells a pint of wine for two asses.}
\end{quote}

Second declension, masculine: \emph{sextārius, sextāriī m.} A sixth of a \emph{congius}, a little over half a litre.

\begin{cousinweb}
\emph{sextārius} is built on \emph{sextus} 'sixth' with the suffix \emph{-ārius} (EtymDB; Lewis \& Short): a sixth of a larger measure.
\end{cousinweb}

\begin{grammarlens}[title={What you've built}]
One more word for telling, explaining and describing at B1.
\end{grammarlens}

\subsection*{Your turn}
\begin{itemize}
\raggedright
  \item \emph{Say it:} \emph{sextārius, sextāriī}
  \item \emph{Say it:} \emph{sextārius, sextāriī}, once more
  \item \emph{Say it:} a sentence of your own with \emph{sextārius, sextāriī}
  \item \emph{Recall:} say the Latin for to weigh carefully, then the Latin for a peck, then say \emph{sextārius, sextāriī} again
\end{itemize}

\begin{tcolorbox}[breakable,colback=teal!4,colframe=teal!35!black,title={Before you move on}]
What does sextārius, sextāriī mean? (\textquotedblleft{}A pint (a sixth of a larger measure)\textquotedblright{}.) Say it once more.
\end{tcolorbox}
\section[pondō]{pondō — by weight, in weight; a pound}

\label{lesson:LA-C465-pondo}



\begin{quote}
Before the new one: say the Latin for a peck, then the Latin for a pint.
\end{quote}

\subsection*{You'll want to know: pondō}
\textbf{pondō} — \textquotedblleft{}by weight, in weight; a pound\textquotedblright{}.

\begin{quote}
Pistor duo pondō farīnae emit.
\end{quote}

\begin{quote}
\emph{The baker buys two pounds of flour.}
\end{quote}

An old ablative of \emph{pondus} 'weight', used as an adverb and as an undeclined noun: \emph{decem pondō}, ten pounds. English \emph{pound} comes from it.

\begin{cousinweb}
\emph{pondō} is an old ablative of \emph{pondus} 'weight' (EtymDB; Lewis \& Short), from the root of \emph{pendō} 'to weigh'. English \emph{pound} comes from it, borrowed early into Germanic.
\end{cousinweb}

\begin{grammarlens}[title={What you've built}]
One more word for telling, explaining and describing at B1.
\end{grammarlens}

\subsection*{Your turn}
\begin{itemize}
\raggedright
  \item \emph{Say it:} \emph{pondō}
  \item \emph{Say it:} \emph{pondō}, once more
  \item \emph{Say it:} a sentence of your own with \emph{pondō}
  \item \emph{Recall:} say the Latin for a peck, then the Latin for a pint, then say \emph{pondō} again
\end{itemize}

\begin{tcolorbox}[breakable,colback=teal!4,colframe=teal!35!black,title={Before you move on}]
What does pondō mean? (\textquotedblleft{}By weight, in weight; a pound\textquotedblright{}.) Say it once more.
\end{tcolorbox}
\section[quadrāns, quadrantis]{quadrāns, quadrantis — a quarter-as (a small coin); a quarter}

\label{lesson:LA-C465-quadrans}



\begin{quote}
Before the new one: say the Latin for a pint, then the Latin for by weight.
\end{quote}

\subsection*{You'll want to know: quadrāns, quadrantis}
\textbf{quadrāns, quadrantis} — \textquotedblleft{}a quarter-as (a small coin); a quarter\textquotedblright{}.

\begin{quote}
Prō balneō tantum quadrantem solvimus.
\end{quote}

\begin{quote}
\emph{We pay only a quarter-as for the bath.}
\end{quote}

Third declension, masculine: \emph{quadrāns, quadrantis m.} A quarter of an \emph{as}, the price of a public bath. English \emph{quadrant} comes from it.

\begin{cousinweb}
\emph{quadrāns} is built on \emph{quattuor} 'four' (EtymDB; Lewis \& Short): a fourth part. English \emph{quadrant} comes from it.
\end{cousinweb}

\begin{grammarlens}[title={What you've built}]
One more word for telling, explaining and describing at B1.
\end{grammarlens}

\subsection*{Your turn}
\begin{itemize}
\raggedright
  \item \emph{Say it:} \emph{quadrāns, quadrantis}
  \item \emph{Say it:} \emph{quadrāns, quadrantis}, once more
  \item \emph{Say it:} a sentence of your own with \emph{quadrāns, quadrantis}
  \item \emph{Recall:} say the Latin for a pint, then the Latin for by weight, then say \emph{quadrāns, quadrantis} again
\end{itemize}

\begin{tcolorbox}[breakable,colback=teal!4,colframe=teal!35!black,title={Before you move on}]
What does quadrāns, quadrantis mean? (\textquotedblleft{}A quarter-as (a small coin); a quarter\textquotedblright{}.) Say it once more.
\end{tcolorbox}
\section[nummulārius, nummulāriī]{nummulārius, nummulāriī — a money-changer, a banker at a table}

\label{lesson:LA-C465-nummularius}



\begin{quote}
Before the new one: say the Latin for by weight, then the Latin for a quarter-as.
\end{quote}

\subsection*{You'll want to know: nummulārius, nummulāriī}
\textbf{nummulārius, nummulāriī} — \textquotedblleft{}a money-changer, a banker at a table\textquotedblright{}.

\begin{quote}
Nummulārius nummōs peregrīnōs in forō mūtat.
\end{quote}

\begin{quote}
\emph{The money-changer changes foreign coins in the forum.}
\end{quote}

Second declension, masculine: \emph{nummulārius, nummulāriī m.} He tested and changed coins at a small table.

\begin{cousinweb}
\emph{nummulārius} is built on \emph{nummulus} 'a little coin, small change', a diminutive of \emph{nummus} 'coin' (EtymDB; Lewis \& Short).
\end{cousinweb}

\begin{grammarlens}[title={What you've built}]
One more word for telling, explaining and describing at B1.
\end{grammarlens}

\subsection*{Your turn}
\begin{itemize}
\raggedright
  \item \emph{Say it:} \emph{nummulārius, nummulāriī}
  \item \emph{Say it:} \emph{nummulārius, nummulāriī}, once more
  \item \emph{Say it:} a sentence of your own with \emph{nummulārius, nummulāriī}
  \item \emph{Recall:} say the Latin for by weight, then the Latin for a quarter-as, then say \emph{nummulārius, nummulāriī} again
\end{itemize}

\begin{tcolorbox}[breakable,colback=teal!4,colframe=teal!35!black,title={Before you move on}]
What does nummulārius, nummulāriī mean? (\textquotedblleft{}A money-changer, a banker at a table\textquotedblright{}.) Say it once more.
\end{tcolorbox}
